% ============ 28 预备约化 (原书 pp.160-162) ============
\chapter{预备约化}\label{sec:28}
\index{弱 Lefschetz 定理}
\index{Gysin 序列}
\index{中间上同调}

我们将证明：只需对偶维数簇的中间上同调群证明定理~\ref{thm:27.15}，而且甚至对这些群，只需证明近似结果.

\begin{lemma}\label{lem:28.1}
只需在把 $\mathbb{F}_q$ 换成 $\mathbb{F}_{q^m}$ 后证明定理~\ref{thm:27.15}.
\end{lemma}

\begin{proof}
相对于 $\mathbb{F}_{q^m}$ 定义的 Frobenius 映射 $F_m \colon X \to X$ 是 $F^{m}$.因此，若 $\alpha_1, \alpha_2, \ldots$ 是 $F$ 在 $H^{r}(X, \mathbb{Q}_{\ell})$ 上的特征值，则 $\alpha^{m}_1, \alpha^{m}_2, \ldots$ 是 $F^{m}$ 在 $H^{r}(X, \mathbb{Q}_{\ell})$ 上的特征值.若 $\alpha^{m}$ 满足相对于 $q^{m}$ 的 27.15 条件，则 $\alpha$ 满足相对于 $q$ 的条件.
\end{proof}

于是，定理~\ref{thm:27.15} 实际上是关于 $X/\mathbb{F}$ 的论断：若它对 $X$ 在有限域上的一个模型成立，则对所有模型都成立.

\begin{exercise}\label{exc:28.2}
设 $X_0$ 是 $\mathbb{F}_q$ 上的三次曲面.用 (\ref{lem:28.1}) 证明 $X_0$ 的 Riemann 猜想.(提示：已知 $X$ 是有理曲面.因此它的 Albanese 簇为零，从而 $H^{1}(X, \mathbb{Q}_{\ell}) = 0$.此外，$H^{2}(X, \mathbb{Q}_{\ell}(1))$ 由 $X$ 上的代数闭链的类生成；事实上，它以 $X$ 上任何 6 条两两不交直线与任何非奇异超平面截的类为基.)
\end{exercise}

\begin{proposition}\label{prop:28.3}
假定对 $\mathbb{F}_q$ 上所有偶维数 $d$ 的非奇异射影簇 $X_0$，$F$ 在 $H^{d}(X, \mathbb{Q}_{\ell})$ 上的每个特征值 $\alpha$ 都是代数数，使得对所有复共轭 $\alpha'$，
\[
q^{\frac{d}{2} - \frac{1}{2}} < |\alpha'| < q^{\frac{d}{2} + \frac{1}{2}}.
\]
则定理~\ref{thm:27.15} 对所有非奇异射影簇成立.
\end{proposition}

\begin{proof}
设 $X_0$ 是 $\mathbb{F}_q$ 上维数 $d$(不一定是偶数)的光滑射影簇，$\alpha$ 是 $F$ 在 $H^{d}(X, \mathbb{Q}_{\ell})$ 上的特征值.K\"unneth 公式表明 $\alpha^{m}$ 出现在 $F$ 在 $H^{dm}(X^{m}, \mathbb{Q}_{\ell})$ 上的特征值中.把引理中的论断应用于 $X_0$ 的偶次幂，表明
\[
q^{\frac{md}{2} - \frac{1}{2}} < |\alpha'|^{m} < q^{\frac{md}{2} + \frac{1}{2}}.
\]
取 $m$ 次根，并让 $m$ 在偶数上趋于 $\infty$，发现
\[
|\alpha'| = q^{\frac{d}{2}}.
\]
现在在命题的假设下对 $X_0$ 的维数作归纳证明 (\ref{thm:27.15}).对 $\dim X_0 = 0$，它显然；对 $d = 1$，只有 $r = 1$ 的情形不显然，而它我们刚刚证明.因此可以假定 $d \geq 2$.

回忆 (\ref{thm:27.12}) 的证明中，Poincar\'e 对偶定理蕴含：若 $\alpha$ 是 $F$ 在 $H^{r}(X, \mathbb{Q}_{\ell})$ 上的特征值，则 $q^{d}/\alpha$ 是 $F$ 在 $H^{2d-r}(X, \mathbb{Q}_{\ell})$ 上的特征值.因此只需对 $r > d$ 证明定理.Bertini 定理(Hartshorne, II.8.18)表明 $\mathbb{P}^{m}$ 中存在超平面 $H$，使得 $Z \overset{\mathrm{def}}{=} H \cap X$ 是非奇异簇.引理~\ref{lem:28.1} 允许我们假定 $H$(从而 $Z$)定义在 $\mathbb{F}_q$ 上.则 Gysin 序列 (\ref{cor:16.2}) 给出
\[
\cdots \to H^{r-2}(Z, \mathbb{Q}_{\ell}(-1)) \to H^{r}(X, \mathbb{Q}_{\ell}) \to H^{r}(X \setminus Z, \mathbb{Q}_{\ell}) \to \cdots.
\]
因为 $X \setminus Z$ 仿射，对 $r > d$，$H^{r}(X \setminus Z, \mathbb{Q}_{\ell}) = 0$(弱 Lefschetz 定理，15.1).于是 Gysin 映射
\[
i_{*} \colon H^{r-2}(Z, \mathbb{Q}_{\ell}(-1)) \to H^{r}(X, \mathbb{Q}_{\ell})
\]
对 $r > d$ 满射.由归纳，$F$ 在 $H^{r-2}(Z, \mathbb{Q}_{\ell})$ 上的特征值是代数数，其共轭有绝对值 $q^{(r-2)/2}$.由于 $F \circ i_{*} = q(i_{*} \circ F)$(\ref{rem:27.4})，$F$ 在 $H^{r}(X, \mathbb{Q}_{\ell})$ 上的特征值是代数数，其共轭有绝对值 $q^{r/2}$.\footnote{来自 Behrooz Mirzaii 与 Steven Spallone 的意见：有可能简化本节的 (\ref{thm:27.15}) 证明，并通过另一个 K\"unneth 公式论证避免使用弱 Lefschetz 定理.我们知道 Frobenius 在 $H^{0}(X)$ 与 $H^{2d}(X)$ 上如何作用.假定 $d < r$.可以把 $H^{r}(X)$ 与自身张量 $d$ 次，把 $H^{0}(X)$ 与自身张量 $r - d$ 次，整个乘积落在 $H^{rd}(X^{rd})$ 中，而我们知道后者.对 $d > r$ 类似.[A. Mellit 的证明.]}
\end{proof}

\subsection*{证明的其余部分}

在 \S 29 中，我们证明仿射曲线上非常值层的 Lefschetz 公式，并在 \S 30 中应用它证明``主引理''，后者给出 $\mathbb{P}^{1}_{/\mathbb{F}_q}$ 的开仿射子集 $U_0$ 上局部常值层 $E$ 的如下判别准则：Frobenius 映射在 $H^{1}(\mathbb{P}^{1}, j_{*}E)$ 上的特征值 $\alpha$ 满足不等式
\[
q^{\frac{n}{2}} < |\alpha| < q^{\frac{n}{2}+1}.
\]
在 \S 31 中，我们研究如何用具有特别好性质(它是 Lefschetz 铅笔)的超曲面截的铅笔 $(X_t)_{t \in \mathbb{P}^{1}}$ 把簇 $X$ 纤维化.在 $X \cap X_t$ 处爆破 $X$ 后，得到簇 $X^{*}$ 与映射 $\varphi \colon X^{*} \to \mathbb{P}^{1}$.在 \S 32 中，我们研究 $\mathbb{Q}_{\ell}$ 在 $\varphi$ 下的高阶直像.

在 \S 33 中，我们把这些主题结合起来完成 Weil 猜想的证明.
